% Options for packages loaded elsewhere
\PassOptionsToPackage{unicode}{hyperref}
\PassOptionsToPackage{hyphens}{url}
\PassOptionsToPackage{dvipsnames,svgnames,x11names}{xcolor}
\documentclass[
  a4paper,
]{article}
\usepackage{xcolor}
\usepackage[margin=25mm]{geometry}
\usepackage{amsmath,amssymb}
\setcounter{secnumdepth}{-\maxdimen} % remove section numbering
\usepackage{iftex}
\ifPDFTeX
  \usepackage[T1]{fontenc}
  \usepackage[utf8]{inputenc}
  \usepackage{textcomp} % provide euro and other symbols
\else % if luatex or xetex
  \usepackage{unicode-math} % this also loads fontspec
  \defaultfontfeatures{Scale=MatchLowercase}
  \defaultfontfeatures[\rmfamily]{Ligatures=TeX,Scale=1}
\fi
\usepackage{lmodern}
\ifPDFTeX\else
  % xetex/luatex font selection
\fi
% Use upquote if available, for straight quotes in verbatim environments
\IfFileExists{upquote.sty}{\usepackage{upquote}}{}
\IfFileExists{microtype.sty}{% use microtype if available
  \usepackage[]{microtype}
  \UseMicrotypeSet[protrusion]{basicmath} % disable protrusion for tt fonts
}{}
\makeatletter
\@ifundefined{KOMAClassName}{% if non-KOMA class
  \IfFileExists{parskip.sty}{%
    \usepackage{parskip}
  }{% else
    \setlength{\parindent}{0pt}
    \setlength{\parskip}{6pt plus 2pt minus 1pt}}
}{% if KOMA class
  \KOMAoptions{parskip=half}}
\makeatother
\usepackage{color}
\usepackage{fancyvrb}
\newcommand{\VerbBar}{|}
\newcommand{\VERB}{\Verb[commandchars=\\\{\}]}
\DefineVerbatimEnvironment{Highlighting}{Verbatim}{commandchars=\\\{\}}
% Add ',fontsize=\small' for more characters per line
\newenvironment{Shaded}{}{}
\newcommand{\AlertTok}[1]{\textcolor[rgb]{1.00,0.00,0.00}{\textbf{#1}}}
\newcommand{\AnnotationTok}[1]{\textcolor[rgb]{0.38,0.63,0.69}{\textbf{\textit{#1}}}}
\newcommand{\AttributeTok}[1]{\textcolor[rgb]{0.49,0.56,0.16}{#1}}
\newcommand{\BaseNTok}[1]{\textcolor[rgb]{0.25,0.63,0.44}{#1}}
\newcommand{\BuiltInTok}[1]{\textcolor[rgb]{0.00,0.50,0.00}{#1}}
\newcommand{\CharTok}[1]{\textcolor[rgb]{0.25,0.44,0.63}{#1}}
\newcommand{\CommentTok}[1]{\textcolor[rgb]{0.38,0.63,0.69}{\textit{#1}}}
\newcommand{\CommentVarTok}[1]{\textcolor[rgb]{0.38,0.63,0.69}{\textbf{\textit{#1}}}}
\newcommand{\ConstantTok}[1]{\textcolor[rgb]{0.53,0.00,0.00}{#1}}
\newcommand{\ControlFlowTok}[1]{\textcolor[rgb]{0.00,0.44,0.13}{\textbf{#1}}}
\newcommand{\DataTypeTok}[1]{\textcolor[rgb]{0.56,0.13,0.00}{#1}}
\newcommand{\DecValTok}[1]{\textcolor[rgb]{0.25,0.63,0.44}{#1}}
\newcommand{\DocumentationTok}[1]{\textcolor[rgb]{0.73,0.13,0.13}{\textit{#1}}}
\newcommand{\ErrorTok}[1]{\textcolor[rgb]{1.00,0.00,0.00}{\textbf{#1}}}
\newcommand{\ExtensionTok}[1]{#1}
\newcommand{\FloatTok}[1]{\textcolor[rgb]{0.25,0.63,0.44}{#1}}
\newcommand{\FunctionTok}[1]{\textcolor[rgb]{0.02,0.16,0.49}{#1}}
\newcommand{\ImportTok}[1]{\textcolor[rgb]{0.00,0.50,0.00}{\textbf{#1}}}
\newcommand{\InformationTok}[1]{\textcolor[rgb]{0.38,0.63,0.69}{\textbf{\textit{#1}}}}
\newcommand{\KeywordTok}[1]{\textcolor[rgb]{0.00,0.44,0.13}{\textbf{#1}}}
\newcommand{\NormalTok}[1]{#1}
\newcommand{\OperatorTok}[1]{\textcolor[rgb]{0.40,0.40,0.40}{#1}}
\newcommand{\OtherTok}[1]{\textcolor[rgb]{0.00,0.44,0.13}{#1}}
\newcommand{\PreprocessorTok}[1]{\textcolor[rgb]{0.74,0.48,0.00}{#1}}
\newcommand{\RegionMarkerTok}[1]{#1}
\newcommand{\SpecialCharTok}[1]{\textcolor[rgb]{0.25,0.44,0.63}{#1}}
\newcommand{\SpecialStringTok}[1]{\textcolor[rgb]{0.73,0.40,0.53}{#1}}
\newcommand{\StringTok}[1]{\textcolor[rgb]{0.25,0.44,0.63}{#1}}
\newcommand{\VariableTok}[1]{\textcolor[rgb]{0.10,0.09,0.49}{#1}}
\newcommand{\VerbatimStringTok}[1]{\textcolor[rgb]{0.25,0.44,0.63}{#1}}
\newcommand{\WarningTok}[1]{\textcolor[rgb]{0.38,0.63,0.69}{\textbf{\textit{#1}}}}
\usepackage{longtable,booktabs,array}
\newcounter{none} % for unnumbered tables
\usepackage{calc} % for calculating minipage widths
% Correct order of tables after \paragraph or \subparagraph
\usepackage{etoolbox}
\makeatletter
\patchcmd\longtable{\par}{\if@noskipsec\mbox{}\fi\par}{}{}
\makeatother
% Allow footnotes in longtable head/foot
\IfFileExists{footnotehyper.sty}{\usepackage{footnotehyper}}{\usepackage{footnote}}
\makesavenoteenv{longtable}
\usepackage{graphicx}
\makeatletter
\newsavebox\pandoc@box
\newcommand*\pandocbounded[1]{% scales image to fit in text height/width
  \sbox\pandoc@box{#1}%
  \Gscale@div\@tempa{\textheight}{\dimexpr\ht\pandoc@box+\dp\pandoc@box\relax}%
  \Gscale@div\@tempb{\linewidth}{\wd\pandoc@box}%
  \ifdim\@tempb\p@<\@tempa\p@\let\@tempa\@tempb\fi% select the smaller of both
  \ifdim\@tempa\p@<\p@\scalebox{\@tempa}{\usebox\pandoc@box}%
  \else\usebox{\pandoc@box}%
  \fi%
}
% Set default figure placement to htbp
\def\fps@figure{htbp}
\makeatother
\setlength{\emergencystretch}{3em} % prevent overfull lines
\providecommand{\tightlist}{%
  \setlength{\itemsep}{0pt}\setlength{\parskip}{0pt}}
\usepackage{bookmark}
\IfFileExists{xurl.sty}{\usepackage{xurl}}{} % add URL line breaks if available
\urlstyle{same}
\hypersetup{
  pdftitle={Mechanism of Inflation-like Rapid Expansion in Self-Consistent Closed Relational-Wave Systems --- Normalization Audit{,} Rank Generation{,} Zero-Square-Closure Conservation{,} Simplex Symmetrization{,} and Reconstruction of the Axiom System},
  pdfauthor={Noriaki Kihara (WF System Co., Ltd.) ORCID 0009-0004-6753-4020},
  colorlinks=true,
  linkcolor={Maroon},
  filecolor={Maroon},
  citecolor={Blue},
  urlcolor={Blue},
  pdfcreator={LaTeX via pandoc}}

\title{Mechanism of Inflation-like Rapid Expansion in Self-Consistent
Closed Relational-Wave Systems --- Normalization Audit, Rank Generation,
Zero-Square-Closure Conservation, Simplex Symmetrization, and
Reconstruction of the Axiom System}
\author{Noriaki Kihara (WF System Co., Ltd.) ORCID 0009-0004-6753-4020}
\date{August 26, 2026 Version DOI 10.5281/zenodo.22112009}

\begin{document}
\maketitle

\textbf{Version DOI:}
\href{https://doi.org/10.5281/zenodo.22112009}{10.5281/zenodo.22112009}\\
\textbf{Concept DOI:}
\href{https://doi.org/10.5281/zenodo.22112008}{10.5281/zenodo.22112008}

\subsection{Abstract}\label{abstract}

Earlier work reported that a closed \(N\)-body relational-wave system,
even with every external seed removed, undergoes a long latency followed
by exponential growth of the component outside the parent plane, passes
through rank-4 and a three-direction readout, and settles into a
metastable structure. This paper reconstructs the onset, the
conservation laws, the saturation, and the geometric ordering of that
process by auditing the code, re-running \(N=3,\ldots,16\), and adding a
linear-stability study at \(N=5\).

First, the Cayley update of the real antisymmetric generator \(K_t\),

\[
C_t=(I-\gamma K_t)^{-1}(I+\gamma K_t),
\]

is a real orthogonal matrix, so \(Z^\dagger Z\) and \(Z^T Z\) are
conserved exactly in exact arithmetic. The zero-square closure that had
been described as ``numerically conserved'' is promoted to a kinematic
theorem; the expansion is not an external injection but an internal
transfer from \(H_\parallel\) to \(H_\perp\), and
\(H_\perp\le H_{\rm total}\) forbids unbounded growth.

Second, when the fixed-point residual of \texttt{make\_parent} at
\(N=5\) is varied from \(3.87\times10^{-7}\) to \(2.38\times10^{-13}\),
the onset step is delayed as 72, 134, 176, 238 while the exponential
growth rate stays at about \(0.17251\)/step. Regression against the
measured fixed-point residual gives

\[
t_{\rm onset}=11.6162[-\ln\varepsilon_{\rm parent}]-99.5631,
\qquad R^2=0.999992 .
\]

Moreover, the rotating-frame Jacobian has a dominant unstable real
doubly degenerate eigenvalue \(\mu_1=1.090086569\) and a weaker second
real doubly degenerate eigenvalue \(\mu_2=1.052603212\). The prediction
from the dominant eigenvalue,

\[
2\ln|\mu_1|=0.172514229,
\]

agrees with the directly measured growth rate of \(H_\perp\), and

\[
1/\ln|\mu_1|=11.5932
\]

agrees with the slope 11.6162 of the onset--residual regression. That
the independent time evolution, the tolerance sweep, and the linearized
spectrum are all consistent with one dominant eigenvalue is the main
quantitative result of this paper. In the numerical spectrum the
dominant multiplier \(\mu_1\) appears as a real doubly degenerate
eigenvalue whose real two-dimensional eigenspace gives the fastest
unstable directions; adding this dominant unstable plane to the rank-2
parent therefore explains the initial selection of rank-4
quantitatively. All twenty multipliers show a clear pair structure, but
fixing the symmetry origin of that doubling analytically would require
showing separately that the Jacobian commutes with the complex
structure.

Third, local zero-square closure at every vertex star is equivalent to
all vertices of the centered complex simplex lying on the bilinear null
cone. Combined with the equimodular amplitudes of the metastable state,
the final geometry is described uniformly as an \textbf{equimodular null
complex simplex}. The complex Gram rank \(N-1\) was confirmed for
\(N=3\)--16.

Fourth, the \(120^\circ\) separation of the three phase classes at
\(N=4\) is derived exactly from local closure and equimodularity. At
\(N=5\), the sign-paired structure \(3+3+2+2\) yields the number of
nontrivial two-edge closures,

\[
3\times3+2\times2=13,
\]

and the number of exact covers partitioning the ten edges into five zero
pairs,

\[
3!\,2!=12,
\]

combinatorially. In an 8-seed sweep, \(3+3+2+2\) and equimodularity were
reproduced in every run, whereas the relative phase between the two
distance families was not fixed uniquely, indicating a flat direction.

Fifth, the axiom system is revised. Complex rotating pairs,
\(\sum z_m^2=0\), and a compact \(S^1\) phase orbit follow from the
self-consistent fixed point, but \(U^n=I\) requires a separate rational
phase-locking mechanism that selects \(\Delta\theta/2\pi\in\mathbb Q\).
The \(K/\sigma_{\max}\) normalization corresponds exactly, at the same
state, to a rescaling of the Cayley step. The roughly 6.8\% difference
in growth rate per accumulated phase observed at finite step is
explained quantitatively by the independent step-refinement series
\(g(n)=g_\infty+c/n\) (\(g_\infty=1.15963\), \(c=-4.105\)): estimating
\(n_{\rm eff}\simeq41.2\) from the mean \(\bar\sigma\simeq3.49\) of the
raw-\(K\) run gives \(g_{\rm pred}=1.0600\), within 0.13\% of the
measured \(1.05874\). Both updates therefore converge to the same
continuum limit, and the possibility that normalization artificially
generated the expansion is excluded both by the Cayley identity and by
the step-convergence series.

\begin{center}\rule{0.5\linewidth}{0.5pt}\end{center}

\subsection{1. Introduction}\label{introduction}

This series studies numerically a closed system in which \(N\) entities
and all their pairwise relations,

\[
M=\frac{N(N-1)}2,
\]

are treated as relational waves. Previously the foundational conditions
were the nontrivial zero-square closure

\[
\sum_m z_m^2=0
\]

and the finite recurrence

\[
U^n=I,
\]

taken as axioms, on top of which the possibilities of finite closed
systems, phase, direction, interaction, and quantization were examined.

The preceding study {[}K8{]} confirmed that even when the external seeds
are removed step by step, the system undergoes geometric rapid expansion
after a long latency, the direction subspace rotates and mixes during
the expansion, and the system finally moves to a metastable
three-direction structure. The simple explanation ``no seed, no
expansion'' was therefore rejected.

The following questions, however, remained open in that study.

\begin{enumerate}
\def\labelenumi{\arabic{enumi}.}
\tightlist
\item
  Why does the rapid expansion start?
\item
  What actually expands?
\item
  When and from where does rank-4 appear?
\item
  Why does the rapid expansion stop?
\item
  Why are directions generated?
\item
  Do the numerical normalization and the step size contained in the
  initialization and the time evolution create the phenomenon?
\item
  Are the two conventional axioms \(\sum z^2=0\) and \(U^n=I\) really
  independent input axioms?
\end{enumerate}

This paper re-examines these questions by an implementation audit and by
re-running \(N=3\)--16.

\begin{center}\rule{0.5\linewidth}{0.5pt}\end{center}

\subsection{2. What Is Inherited from the Preceding Study and What Is
Corrected}\label{what-is-inherited-from-the-preceding-study-and-what-is-corrected}

\subsubsection{2.1 The direct predecessor}\label{the-direct-predecessor}

The direct predecessor is {[}K8{]}. It confirmed that rapid expansion
and direction generation persist after seed removal.

That paper, however, contained the following statement.

\begin{quote}
The initialization \texttt{make\_parent} presupposes the zero-square
closure condition.
\end{quote}

Auditing the code back to the original source, we found that this
statement needs correction.

\texttt{make\_parent} does not impose \(\sum z_m^2=0\) directly as a
numerical constraint. What it actually does is search for a
self-consistent eigenmode fixed point of a phase-dependent real
antisymmetric action. Zero-square closure appears on the output side as
a consequence.

This paper therefore does not deny the experimental facts of {[}K8{]};
it \textbf{corrects the understanding of the initialization mechanism
through a code audit and, as a consequence, reconstructs the axiom
system further upstream}.

\subsubsection{2.2 What was re-examined}\label{what-was-re-examined}

We re-audited

\begin{itemize}
\tightlist
\item
  the original code,
\item
  \texttt{make\_parent},
\item
  the Cayley time evolution,
\item
  the \(K/\sigma_{\max}\) normalization,
\item
  the step parameter \texttt{GAMMA},
\item
  the time evolution of the rank,
\item
  \(Z^T Z\),
\item
  complex distances,
\item
  the simplex rank,
\item
  the metastabilization of amplitudes and phases,
\item
  partial zero closures.
\end{itemize}

\begin{center}\rule{0.5\linewidth}{0.5pt}\end{center}

\subsection{3. Minimal Input of the
System}\label{minimal-input-of-the-system}

The only independent size parameter in this paper is \(N\).

Because all pairwise relations of \(N\) entities are taken,

\[
M=\binom N2=\frac{N(N-1)}2
\]

is a derived quantity, not an independent parameter.

Abstracting the current implementation, the minimal input is

\begin{center}\fbox{\parbox{0.9\linewidth}{\centering $N
+$ complete pairwise relations $+$ self-consistent fixed-point condition}}\end{center}

\begin{center}\rule{0.5\linewidth}{0.5pt}\end{center}

\subsection{\texorpdfstring{4. Implementation Audit of
\texttt{make\_parent}}{4. Implementation Audit of make\_parent}}\label{implementation-audit-of-make_parent}

\subsubsection{4.1 The actual code}\label{the-actual-code}

The core of the \texttt{make\_parent} currently in use is the following.

\begin{Shaded}
\begin{Highlighting}[]
\KeywordTok{def}\NormalTok{ make\_parent(sys\_lr, rng, iters}\OperatorTok{=}\DecValTok{400}\NormalTok{, beta}\OperatorTok{=}\FloatTok{0.5}\NormalTok{, tol}\OperatorTok{=}\FloatTok{1e{-}8}\NormalTok{, restarts}\OperatorTok{=}\DecValTok{3}\NormalTok{):}
\NormalTok{    best }\OperatorTok{=}\NormalTok{ (}\VariableTok{None}\NormalTok{, np.inf, }\VariableTok{None}\NormalTok{)}
    \ControlFlowTok{for}\NormalTok{ \_ }\KeywordTok{in} \BuiltInTok{range}\NormalTok{(restarts):}
\NormalTok{        theta }\OperatorTok{=}\NormalTok{ rng.uniform(}\FloatTok{0.0}\NormalTok{, }\FloatTok{2.0} \OperatorTok{*}\NormalTok{ np.pi, sys\_lr.m)}
\NormalTok{        v }\OperatorTok{=} \VariableTok{None}
        \ControlFlowTok{for}\NormalTok{ it }\KeywordTok{in} \BuiltInTok{range}\NormalTok{(iters):}
\NormalTok{            sys\_lr.set\_theta(theta)}
\NormalTok{            ev, EV }\OperatorTok{=}\NormalTok{ np.linalg.eig(sys\_lr.J }\OperatorTok{@}\NormalTok{ sys\_lr.G)}
\NormalTok{            idx }\OperatorTok{=} \BuiltInTok{int}\NormalTok{(np.argmin(ev.imag))  }\CommentTok{\# lambda = {-}i sigma\_max}
\NormalTok{            v }\OperatorTok{=}\NormalTok{ sys\_lr.w(EV[:, idx].astype(}\BuiltInTok{complex}\NormalTok{))}
\NormalTok{            v }\OperatorTok{=}\NormalTok{ v }\OperatorTok{/}\NormalTok{ np.linalg.norm(v)}
\NormalTok{            theta\_new }\OperatorTok{=}\NormalTok{ np.angle(v)}
\NormalTok{            mix }\OperatorTok{=}\NormalTok{ (}\FloatTok{1.0} \OperatorTok{{-}}\NormalTok{ beta) }\OperatorTok{*}\NormalTok{ np.exp(}\OtherTok{1j} \OperatorTok{*}\NormalTok{ theta) }\OperatorTok{\textbackslash{}}
                \OperatorTok{+}\NormalTok{ beta }\OperatorTok{*}\NormalTok{ np.exp(}\OtherTok{1j} \OperatorTok{*}\NormalTok{ theta\_new)}
\NormalTok{            theta }\OperatorTok{=}\NormalTok{ np.angle(mix)}

            \ControlFlowTok{if}\NormalTok{ it }\OperatorTok{\%} \DecValTok{10} \OperatorTok{==} \DecValTok{9}\NormalTok{:}
\NormalTok{                sys\_lr.set\_theta(np.angle(v))}
\NormalTok{                res\_now }\OperatorTok{=}\NormalTok{ \_eigenmode\_residual(sys\_lr, v)}
                \ControlFlowTok{if}\NormalTok{ res\_now }\OperatorTok{\textless{}}\NormalTok{ tol:}
                    \ControlFlowTok{break}

\NormalTok{        sys\_lr.set\_theta(np.angle(v))}
\NormalTok{        residual }\OperatorTok{=}\NormalTok{ \_eigenmode\_residual(sys\_lr, v)}
        \ControlFlowTok{if}\NormalTok{ residual }\OperatorTok{\textless{}}\NormalTok{ best[}\DecValTok{1}\NormalTok{]:}
\NormalTok{            best }\OperatorTok{=}\NormalTok{ (v, residual, sys\_lr.sigma\_spectrum())}
        \ControlFlowTok{if}\NormalTok{ residual }\OperatorTok{\textless{}}\NormalTok{ tol:}
            \ControlFlowTok{break}

\NormalTok{    v, residual, sig }\OperatorTok{=}\NormalTok{ best}
\NormalTok{    sys\_lr.set\_theta(np.angle(v))}
    \ControlFlowTok{return}\NormalTok{ v, residual, sig}
\end{Highlighting}
\end{Shaded}

Here the initial value \texttt{theta} is

\begin{Shaded}
\begin{Highlighting}[]
\NormalTok{theta }\OperatorTok{=}\NormalTok{ rng.uniform(}\FloatTok{0.0}\NormalTok{, }\FloatTok{2.0}\OperatorTok{*}\NormalTok{np.pi, M)}
\end{Highlighting}
\end{Shaded}

that is, \textbf{real random numbers}.

No pair of random amplitude sets for real and imaginary parts is given
as a physical initial value.

The \texttt{astype(complex)} in the implementation exists only to store
the complex-conjugate eigenmode of the real matrix \texttt{J\ @\ G}; it
does not specify any complex phase relation as an external condition.

\subsubsection{4.2 Self-consistency as a fixed
point}\label{self-consistency-as-a-fixed-point}

Conceptually, the iteration

\[
\theta
\rightarrow
G(\theta)
\rightarrow
v
\rightarrow
\arg v
\rightarrow
\theta'
\]

is repeated and a fixed point with

\[
\boxed{\theta'=\theta}
\]

is sought.

The essence of \texttt{make\_parent} is therefore the self-consistent
fixed-point condition

\[
\boxed{\mathcal F(X)=X} .
\]

\begin{center}\rule{0.5\linewidth}{0.5pt}\end{center}

\subsection{5. Spontaneous Generation of the Complex
Structure}\label{spontaneous-generation-of-the-complex-structure}

\subsubsection{5.1 Real antisymmetric
action}\label{real-antisymmetric-action}

The generator \(K\) is real antisymmetric,

\[
K^T=-K.
\]

Its nonzero eigenmodes have the form

\[
Kv=i\sigma v .
\]

Writing

\[
v=a+ib
\]

gives

\[
Ka=-\sigma b,\qquad
Kb=\sigma a.
\]

Thus \(a,b\) are not two components supplied independently from outside;
they are \textbf{the rotating pair that appears inside a nonzero
eigenmode of a real antisymmetric action}.

In this sense one may read

\begin{center}\fbox{\parbox{0.9\linewidth}{\centering real initial phases $+$ self-consistent real antisymmetric action $\Rightarrow$ complex rotating structure}}\end{center}

\begin{center}\rule{0.5\linewidth}{0.5pt}\end{center}

\subsection{6. From Self-Consistency to Zero-Square
Closure}\label{from-self-consistency-to-zero-square-closure}

By real antisymmetry, for any \(v\),

\[
v^T K v=0 .
\]

Substituting a self-consistent nonzero eigenmode

\[
Kv=i\sigma v,\qquad \sigma\neq0
\]

gives

\[
v^TKv=i\sigma v^Tv=0 .
\]

Hence

\[
\boxed{v^Tv=0},
\]

that is,

\[
\boxed{\sum_m v_m^2=0} .
\]

Furthermore, with \(v=a+ib\),

\[
\sum_m(a_m+ib_m)^2=0
\]

yields

\[
\boxed{\|a\|^2=\|b\|^2},
\qquad
\boxed{a\cdot b=0}.
\]

Zero-square closure is therefore not an external constraint given to
\texttt{make\_parent}; it is derived from the self-consistent nonzero
antisymmetric eigenmode.

\begin{center}\rule{0.5\linewidth}{0.5pt}\end{center}

\subsection{7. Linear-Combination Closure at Each
Vertex}\label{linear-combination-closure-at-each-vertex}

In the complete graph on \(N\) vertices, \(N-1\) relational waves are
incident to vertex \(i\).

In the present \(N=3\)--16 analysis, for \(N\ge4\) each vertex satisfied

\[
\boxed{
\sum_{j\neq i} z_{ij}^2\simeq0
}
\]

within numerical precision.

This means that the state of each vertex is a linear-combination closure
built from the \(N-1\) relational waves other than itself:

\begin{center}\fbox{\parbox{0.9\linewidth}{\centering state of one vertex $=$ self-consistent linear combination of the $(N-1)$ incident relational waves}}\end{center}

Since each relation \(z_{ij}\) belongs to the local closures of both
vertex \(i\) and vertex \(j\), the whole system is covered by \(N\)
mutually overlapping local linear closures.

\begin{center}\rule{0.5\linewidth}{0.5pt}\end{center}

\subsection{\texorpdfstring{8. From Zero-Square Closure to a Compact
Phase Orbit --- the Logical Distinction from
\(U^n=I\)}{8. From Zero-Square Closure to a Compact Phase Orbit --- the Logical Distinction from U\^{}n=I}}\label{from-zero-square-closure-to-a-compact-phase-orbit-the-logical-distinction-from-uni}

From

\[
v^Tv=0
\]

we have, for \(v=a+ib\),

\[
\|a\|=\|b\|,
\qquad
a\perp b,
\]

so the state can be written as a rotating pair of finite radius,

\[
X(\theta)=a\cos\theta+b\sin\theta .
\]

Hence

\[
X(\theta+2\pi)=X(\theta),
\]

and a \textbf{compact \(S^1\) phase orbit} can be derived from the
self-consistent fixed point.

Finite recurrence,

\[
U^n=I,
\]

however, does not follow automatically. To return to the initial state
after a finite \(n\), the phase increment must satisfy

\[
\frac{\Delta\theta}{2\pi}\in\mathbb Q .
\]

For an irrational rotation the orbit is dense on \(S^1\) and never
closes in finitely many steps.

The logical hierarchy that this paper can establish rigorously is
therefore

\begin{center}\fbox{\parbox{0.9\linewidth}{\centering self-consistent fixed point $\Rightarrow$ complex rotating pair $\Rightarrow
\sum z_m^2=0
\Rightarrow
S^1$ compactness}}\end{center}

and no further. The remaining step must be separated as

\begin{center}\fbox{\parbox{0.9\linewidth}{\centering $S^1$ compactness $+$ rational phase locking / discrete phase selection $\Rightarrow
U^n=I$}}\end{center}

This paper therefore does not claim that \(U^n=I\), previously placed as
an axiom, has been fully derived from self-consistency. The rational
phase-locking mechanism is an independent problem to be examined in
connection with the separate phase-discretization experiments.

\subsection{9. Implementation of the Time
Evolution}\label{implementation-of-the-time-evolution}

At each step the time evolution reads the phase of the current state,

\begin{Shaded}
\begin{Highlighting}[]
\NormalTok{sys\_lr.set\_theta(np.angle(Z))}
\NormalTok{sig\_est, wp }\OperatorTok{=}\NormalTok{ sys\_lr.sigma\_max\_power(wp)}
\NormalTok{Z }\OperatorTok{=}\NormalTok{ sys\_lr.cayley\_step(Z, sig\_est)}
\end{Highlighting}
\end{Shaded}

and writes back the next state.

Physically,

\[
Z_t
\rightarrow
\theta_t=\arg Z_t
\rightarrow
K(\theta_t)
\rightarrow
U_t Z_t
\rightarrow
Z_{t+1} .
\]

The generator is therefore not fixed: it is \textbf{a state-dependent
generator reconstructed at every step from the readout of the current
wave state}.

\begin{center}\rule{0.5\linewidth}{0.5pt}\end{center}

\subsection{10. Cayley Update and Phase
Rotation}\label{cayley-update-and-phase-rotation}

The time evolution is the Cayley transform

\[
Z_{t+1}
=
(I-\gamma K_t)^{-1}
(I+\gamma K_t)Z_t .
\]

For real antisymmetric \(K_t\) this update is a norm-preserving
rotation.

For an eigenmode

\[
K_tv=i\sigma v
\]

we have

\[
v\rightarrow
\frac{1+i\gamma\sigma}
     {1-i\gamma\sigma}v
=
e^{i\Delta\phi}v,
\]

\[
\boxed{
\Delta\phi
=
2\arctan(\gamma\sigma)
} .
\]

This formula is the mathematical expression of ``read the state, rotate
the phase, write back the next state''.

\begin{center}\rule{0.5\linewidth}{0.5pt}\end{center}

\subsection{\texorpdfstring{11. The Unnecessary \(K/\sigma_{\max}\)
Normalization and Time
Reparametrization}{11. The Unnecessary K/\textbackslash sigma\_\{\textbackslash max\} Normalization and Time Reparametrization}}\label{the-unnecessary-ksigma_max-normalization-and-time-reparametrization}

\subsubsection{11.1 Old code and corrected
code}\label{old-code-and-corrected-code}

The earlier code used the generator

\[
K\rightarrow\frac{K}{\sigma_{\max}} .
\]

The corrected code uses the raw \(K\).

At the same state the Cayley map satisfies the identity

\[
C\!\left(\frac K\sigma,\gamma\right)
=
C\!\left(K,\frac\gamma\sigma\right) .
\]

In the continuum limit, therefore, the vector fields with and without
normalization point in the same direction, and the normalization can be
understood as a state-dependent change of clock.

In the earlier step-axis comparison at \(N=4,5\), the ratio of growth
rates raw/normalized and the ratio of arrival steps normalized/raw were

\[
N=4:\quad2.708\ \text{vs}\ 2.702,
\]

\[
N=5:\quad3.495\ \text{vs}\ 3.489,
\]

in agreement, which strongly suggested the time-reparametrization
interpretation.

\subsubsection{11.2 Additional verification on the accumulated-phase
axis}\label{additional-verification-on-the-accumulated-phase-axis}

The existing \(N=5\) data were re-plotted against the accumulated Cayley
phase of the dominant mode,

\[
\Phi(t)=\sum_{s<t}2\arctan(\gamma\sigma_s) .
\]

\begin{figure}
\centering
\pandocbounded{\includegraphics[keepaspectratio,alt={N=5: comparison with and without normalization on the accumulated-phase axis. The finite-step trajectory difference is compared in terms of the accumulated Cayley phase.}]{figs/N5_normalized_vs_raw_cumulative_phase.png}}
\caption{\(N=5\): comparison with and without normalization on the
accumulated-phase axis. The finite-step trajectory difference is
compared in terms of the accumulated Cayley phase.}
\end{figure}

At finite step the two curves do not coincide completely; the RMSE of
\(\log_{10}H_\perp\) in the growth region was \textbf{0.944 decade}. In
an unstable system, however, a tiny discretization difference is
amplified exponentially, so the trajectory RMSE alone is not used as a
refutation of the time-reparametrization hypothesis. We therefore
regressed separately the slope of \(\ln H_\perp\) per accumulated phase
in the exponential-growth region \(10^{-10}<H_\perp<10^{-3}\) of the
same data. The result is

\[
\frac{d\ln H_\perp}{d\Phi}
=
1.13117\quad(K/\sigma_{\max}),
\]

\[
\frac{d\ln H_\perp}{d\Phi}
=
1.05874\quad(\text{raw }K),
\]

a difference of about 6.8\%. The intercepts of the same regressions,
\(-55.7414\) and \(-55.7412\), almost coincide. At the current
\(\gamma=\tan(\pi/144)\), therefore, the difference cannot be absorbed
by a mere horizontal or vertical offset; a small but measurable
difference remains in the effective finite-step growth rate.

This 6.8\% difference can be closed with the step-refinement series
already obtained in the next section. A least-squares fit of the five
points of §12 to

\[
g(n)=g_\infty+\frac{c}{n}
\]

gives

\[
g_\infty=1.1596346,\qquad c=-4.10498,
\]

with a maximum residual of \(8.12\times10^{-5}\) over the five points.
The growth rate per accumulated phase is thus described to high
precision by a first-order finite-step correction.

On the other hand, by the identity at the same state,

\[
C(K/\sigma,\gamma)=C(K,\gamma/\sigma),
\]

the raw-\(K\) update can be read as the normalized update run at a
coarser effective step on average. Estimating \(\bar\sigma\simeq3.49\)
from the \(N=5\) step-axis growth-rate ratio,

\[
n_{\rm eff}\simeq\frac{144}{3.49}=41.2 .
\]

At this point the independent convergence law predicts

\[
g(41.2)=1.1596346-\frac{4.10498}{41.2}\simeq1.0600,
\]

which agrees with the directly measured raw-\(K\) value \(1.05874\) to
about 0.13\%. Since \(\sigma_t\) is time dependent, \(n_{\rm eff}\) is
an average approximation, but this agreement shows that the roughly
6.8\% finite-step difference is explained by the same first-order step
correction. The agreement of the two intercepts, \(-55.7414\) and
\(-55.7412\), is also consistent with the interpretation that the
initial perturbation is common and the difference stems mainly from the
step.

The conclusion is therefore

\begin{center}\fbox{\parbox{0.9\linewidth}{\centering the $K/\sigma_{\max}$ normalization is a state-dependent rescaling of the Cayley step; the finite-step raw/normalized difference is explained quantitatively by the independently measured first-order step correction, and both converge to the same continuum limit.}}\end{center}

Hence the possibility that normalization or step size artificially
generated the inflation-like expansion itself is excluded both by the
algebraic identity and by the continuum-limit series.

\subsection{\texorpdfstring{12. \texttt{144} Is a Time Step, Not a
Physical
Constant}{12. 144 Is a Time Step, Not a Physical Constant}}\label{is-a-time-step-not-a-physical-constant}

The code sets

\begin{Shaded}
\begin{Highlighting}[]
\NormalTok{GAMMA }\OperatorTok{=}\NormalTok{ math.tan(math.pi }\OperatorTok{/} \FloatTok{144.0}\NormalTok{)}
\end{Highlighting}
\end{Shaded}

Note that \(\pi/144\) is the angle of the Cayley parameter; the actual
one-step rotation angle for a normalized eigenmode with \(\sigma=1\) is

\[
\Delta\phi
=
2\arctan[\tan(\pi/144)]
=
\frac{2\pi}{144}
=
\frac{\pi}{72}
=
2.5^\circ .
\]

Thus ``144'' is not a physical constant but the denominator of the
numerical time step chosen to obtain a \(2.5^\circ\) increment.

For \(N=5\) we refined the step as

\[
n_\gamma=144,\ 288,\ 576,\ 1152,\ 2304
\]

with

\[
\gamma=\tan(\pi/n_\gamma)
\]

and examined the continuum limit.

The growth rate per accumulated phase converged as

\begin{itemize}
\tightlist
\item
  144: 1.131171 /rad
\item
  288: 1.145300 /rad
\item
  576: 1.152475 /rad
\item
  1152: 1.156090 /rad
\item
  2304: 1.157905 /rad
\end{itemize}

These five points are described by the first-order convergence law

\[
g(n)=g_\infty+\frac{c}{n},\qquad
g_\infty=1.1596346,\quad c=-4.10498,
\]

with a maximum residual of \(8.12\times10^{-5}\). This convergence law
independently explains the raw/normalized difference of §11.2.

Therefore

\begin{center}\fbox{\parbox{0.9\linewidth}{\centering there is no physical resonance specific to 144; the dynamics converges to the continuum-time limit $\gamma\rightarrow0 .$}}\end{center}

\begin{figure}
\centering
\pandocbounded{\includegraphics[keepaspectratio,alt={N=5: continuum limit under step refinement. The growth rate per accumulated phase converges to a common limit with a first-order finite-step correction.}]{figs/N5_gamma_continuum_convergence.png}}
\caption{\(N=5\): continuum limit under step refinement. The growth rate
per accumulated phase converges to a common limit with a first-order
finite-step correction.}
\end{figure}

\begin{center}\rule{0.5\linewidth}{0.5pt}\end{center}

\subsection{13. The Rank-2 Initial State and the Spontaneous Generation
of
Rank-4}\label{the-rank-2-initial-state-and-the-spontaneous-generation-of-rank-4}

The parent state obtained by \texttt{make\_parent} has a rank-2 initial
plane spanned by its real and imaginary parts.

When the time evolution starts from this state, the component outside
the parent plane rises spontaneously and the readout structure moves to
rank-4.

The important point is

\[
\boxed{
\text{rank-4 is not given as an initial condition}
} .
\]

Rank-4 is generated spontaneously by the iteration that reads the
current state, reconstructs the generator, and writes back the
phase-rotated result.

The rank-4 referred to here differs from the Gram rank \(N-1\) of the
complex simplex discussed later. The former is the readout of the ``two
two-dimensional rotation planes'' that rise during the time evolution;
the latter is the simplex rank of the distance geometry of \(N\)
vertices.

\begin{center}\rule{0.5\linewidth}{0.5pt}\end{center}

\subsection{14. Generation of Three
Directions}\label{generation-of-three-directions}

Let the initial rank-2 plane be \(A\).

During the time evolution a second two-dimensional plane \(B\) rises
outside the parent plane.

Reading the relative configuration of the two planes, one can identify
three directions:

\begin{enumerate}
\def\labelenumi{\arabic{enumi}.}
\tightlist
\item
  the shared direction related to both planes,
\item
  the independent direction on the side of plane \(A\),
\item
  the independent direction on the side of plane \(B\).
\end{enumerate}

Hence the generation sequence

\begin{center}\fbox{\parbox{0.9\linewidth}{\centering rank-2 $\rightarrow$ rank-4 $\rightarrow$ three-direction readout}}\end{center}

is obtained.

The ``three directions'' observed in {[}K8{]} are not an amplification
of three axes fixed from the beginning; they are formed dynamically
during the rapid expansion, accompanied by rotation and mixing of
subspaces.

\begin{center}\rule{0.5\linewidth}{0.5pt}\end{center}

\subsection{15. What Expands --- Internal Transfer under Exact
Conservation}\label{what-expands-internal-transfer-under-exact-conservation}

Let \(Z_\parallel,Z_\perp\) be the projection onto the parent plane and
the orthogonal component, and

\[
H_\parallel=\|Z_\parallel\|^2,
\qquad
H_\perp=\|Z_\perp\|^2 .
\]

Since the orthogonal decomposition is fixed, together with the exact
conservation law of the Cayley update shown in the next section,

\[
H_\parallel+H_\perp=H_{\rm total} .
\]

We overlaid both quantities directly from the corrected raw-\(K\)
\(N=5\) data.

\begin{figure}
\centering
\pandocbounded{\includegraphics[keepaspectratio,alt={N=5: depletion of the parent-plane component and growth of the transverse component. Their sum is conserved exactly, showing that the rapid expansion is an internal transfer.}]{figs/N5_pump_depletion_Hparallel_Hperp.png}}
\caption{\(N=5\): depletion of the parent-plane component and growth of
the transverse component. Their sum is conserved exactly, showing that
the rapid expansion is an internal transfer.}
\end{figure}

The maximum error of the implementation is

\[
\max|H_\parallel+H_\perp-H_{\rm total}|=4.44\times10^{-16}.
\]

At step 5000,

\[
H_\parallel=0.3220054054,
\qquad
H_\perp=0.6779945946 .
\]

The rapid growth of \(H_\perp\) is therefore not an external injection
but an \textbf{exact internal redistribution} from the parent-plane
component to the orthogonal directions.

Moreover, since

\[
0\le H_\perp\le H_{\rm total}=\text{const.},
\]

the fact that the exponential growth does not continue indefinitely is
kinematically enforced.

\subsection{16. Two Exact Conservation Laws of the Cayley
Update}\label{two-exact-conservation-laws-of-the-cayley-update}

The generator at each step is real antisymmetric,

\[
K_t^T=-K_t .
\]

Hence the Cayley map

\[
C_t=(I-\gamma K_t)^{-1}(I+\gamma K_t)
\]

is a real orthogonal matrix satisfying identically

\[
C_t^T C_t=I .
\]

In exact arithmetic, therefore, even though the generator changes with
the state at every step,

\[
\boxed{Z_{t+1}^\dagger Z_{t+1}=Z_t^\dagger Z_t}
\]

and

\[
\boxed{Z_{t+1}^T Z_{t+1}=Z_t^T Z_t}
\]

are conserved exactly.

The earlier statements ``\(H_{\rm total}\) is almost constant'' and
``\(Z^TZ\simeq0\) is maintained'' are thus promoted from numerical
findings to \textbf{structural theorems of the update rule}. The
residuals of order \(10^{-15}\) observed in the \(N=5\) raw-\(K\)
implementation are the floating-point implementation error of this exact
theorem.

This theorem guarantees analytically that the inflation-like expansion
is not an energy generation that violates conservation, but an internal
rearrangement of a conserved total quantity.

\subsection{17. Complex Simplex and the Null Cone --- a Geometric
Theorem on Local
Closure}\label{complex-simplex-and-the-null-cone-a-geometric-theorem-on-local-closure}

Let the complex squared distance of each relation be

\[
q_{ij}=z_{ij}^2=(x_i-x_j)\cdot(x_i-x_j),
\]

and take the centroid at

\[
\sum_i x_i=0 .
\]

The star sum at vertex \(i\) is

\[
\sum_{j\ne i}(x_i-x_j)^2
=Nx_i^2+T,
\qquad
T\equiv\sum_jx_j^2 .
\]

If local closure

\[
\sum_{j\ne i}q_{ij}=0
\]

holds at every vertex, then

\[
x_i^2=-T/N
\]

is common to all \(i\). Summing over all vertices gives \(T=-T\), hence

\[
T=0,
\]

and therefore

\[
\boxed{x_i\cdot x_i=0\quad\forall i}.
\]

The converse holds immediately, so

\begin{center}\fbox{\parbox{0.9\linewidth}{\centering zero-square closure of every vertex star $\Longleftrightarrow$ all simplex vertices lie on the complex null cone}}\end{center}

is a theorem.

In this case

\[
q_{ij}=-2x_i\cdot x_j .
\]

Across the whole series \(N=3\)--16 the complex Gram rank is

\[
\operatorname{rank}B=N-1,
\]

so the \(N\) vertices form an \((N-1)\)-dimensional complex simplex.

\subsection{18. Metastabilization to an Equimodular Null Complex
Simplex}\label{metastabilization-to-an-equimodular-null-complex-simplex}

After long evolution, for most \(N\),

\[
|z_{ij}|^2\rightarrow\frac1M ,
\]

i.e.~equipartition.

Combined with the null-cone theorem of the previous section, the
metastable state is summarized as one geometric object,

\[
\boxed{
\text{equimodular null complex simplex}
} ,
\]

namely

\[
x_i^2=0,
\qquad
|x_i\cdot x_j|=\frac{1}{2M}\quad(i\ne j).
\]

For \(N=5\) we computed the spectral entropy of the amplitude
distribution,

\[
p_m=\frac{|z_m|^2}{\sum_k|z_k|^2},
\qquad
S=-\sum_mp_m\ln p_m .
\]

\begin{figure}
\centering
\pandocbounded{\includegraphics[keepaspectratio,alt={N=5: spectral entropy. It approaches equipartition over long times, but not as a simple monotone H-theorem.}]{figs/N5_spectral_entropy.png}}
\caption{\(N=5\): spectral entropy. It approaches equipartition over
long times, but not as a simple monotone H-theorem.}
\end{figure}

Initially \(S/\ln M=0.97472\); it dipped once to 0.97322 at step 375 and
reached

\[
\boxed{S/\ln M=1.000000}
\]

at step 5000. The final equipartition is therefore quantitatively
complete. Since \(S(t)\) is not strictly monotone, however, this differs
from a simple H-theorem.

\subsection{19. Stopping of the Rapid Expansion: Kinematic Bound and
Subsequent Simplex
Ordering}\label{stopping-of-the-rapid-expansion-kinematic-bound-and-subsequent-simplex-ordering}

By the Cayley orthogonality,

\[
H_\perp\le H_{\rm total}=\text{const.},
\]

so the fact that the exponential growth saturates rather than continuing
indefinitely is explained kinematically.

The question ``why does it stop'' therefore splits into two.

\begin{enumerate}
\def\labelenumi{\arabic{enumi}.}
\tightlist
\item
  \textbf{Why it does not grow indefinitely}: resolved by the bound from
  total-norm conservation.
\item
  \textbf{Why it moves to a specific equimodular simplex configuration}:
  still a dynamical problem.
\end{enumerate}

At \(N=5\), even after \(H_\perp\) has nearly saturated, the ordering of
the four phase/distance classes continues for a long time.

\begin{figure}
\centering
\pandocbounded{\includegraphics[keepaspectratio,alt={N=5: time evolution of rapid expansion and ordering. Simplex ordering continues after the expansion has saturated.}]{figs/N5_inflation_vs_ordering.png}}
\caption{\(N=5\): time evolution of rapid expansion and ordering.
Simplex ordering continues after the expansion has saturated.}
\end{figure}

The time evolution therefore has at least two stages,

\begin{center}\fbox{\parbox{0.9\linewidth}{\centering linear instability / rapid transfer $\rightarrow$ bounded saturation $\rightarrow$ slow simplex ordering}}\end{center}

\subsection{\texorpdfstring{20. The Series
\(N=3\)--16}{20. The Series N=3--16}}\label{the-series-n316}

The main results are summarized below.

{\def\LTcaptype{none} % do not increment counter
\begin{longtable}[]{@{}
  >{\raggedleft\arraybackslash}p{(\linewidth - 6\tabcolsep) * \real{0.2667}}
  >{\raggedleft\arraybackslash}p{(\linewidth - 6\tabcolsep) * \real{0.2667}}
  >{\raggedleft\arraybackslash}p{(\linewidth - 6\tabcolsep) * \real{0.2667}}
  >{\raggedright\arraybackslash}p{(\linewidth - 6\tabcolsep) * \real{0.2000}}@{}}
\toprule\noalign{}
\begin{minipage}[b]{\linewidth}\raggedleft
N
\end{minipage} & \begin{minipage}[b]{\linewidth}\raggedleft
M
\end{minipage} & \begin{minipage}[b]{\linewidth}\raggedleft
simplex rank
\end{minipage} & \begin{minipage}[b]{\linewidth}\raggedright
outline of the metastable phase structure
\end{minipage} \\
\midrule\noalign{}
\endhead
\bottomrule\noalign{}
\endlastfoot
3 & 3 & 2 & triangular simplex, still relaxing at 5000 steps \\
4 & 6 & 3 & \(2+2+2\): the three pairs of opposite edges of a
tetrahedron \\
5 & 10 & 4 & \(3+3+2+2\), plus nontrivial two-edge closures \\
6 & 15 & 5 & strictly multi-class, a boundary case with slow
relaxation \\
7 & 21 & 6 & 21 edges nondegenerate at high precision \\
8 & 28 & 7 & 28 edges nondegenerate at high precision \\
9 & 36 & 8 & 36 edges nondegenerate at high precision \\
10 & 45 & 9 & 45 edges nondegenerate at high precision \\
11 & 55 & 10 & 55 edges nondegenerate at high precision \\
12 & 66 & 11 & 66 edges nondegenerate at high precision \\
13 & 78 & 12 & 78 edges nondegenerate at high precision \\
14 & 91 & 13 & nondegenerate, one 6-edge quasi-closure candidate \\
15 & 105 & 14 & 105 edges nondegenerate at high precision \\
16 & 120 & 15 & 120 edges nondegenerate at high precision \\
\end{longtable}
}

Universal are

\[
\boxed{
\operatorname{rank}=N-1
}
\]

and, in the metastable state,

\[
\boxed{
|z|^2\rightarrow1/M
} .
\]

The small phase groups and the additional closures, on the other hand,
depend strongly on \(N\).

\begin{center}\rule{0.5\linewidth}{0.5pt}\end{center}

\subsection{\texorpdfstring{21. The Opposite-Edge Structure at \(N=4\)
--- \(120^\circ\) Is a Theorem, Not a Numerical
Approximation}{21. The Opposite-Edge Structure at N=4 --- 120\^{}\textbackslash circ Is a Theorem, Not a Numerical Approximation}}\label{the-opposite-edge-structure-at-n4-120circ-is-a-theorem-not-a-numerical-approximation}

At \(N=4\) the six edges split into three pairs of opposite edges,

\[
(12,34),\qquad(13,24),\qquad(14,23).
\]

\begin{figure}
\centering
\pandocbounded{\includegraphics[keepaspectratio,alt={N=4: the three opposite-edge classes of the tetrahedron. The 120\^{}\textbackslash circ separation of the three classes follows exactly from local closure and equimodularity.}]{figs/N4_tetrahedron_opposite_edge_classes.png}}
\caption{\(N=4\): the three opposite-edge classes of the tetrahedron.
The \(120^\circ\) separation of the three classes follows exactly from
local closure and equimodularity.}
\end{figure}

Each vertex star contains one edge from each class, so local closure
reads

\[
v_1+v_2+v_3=0 .
\]

Under metastable equipartition,

\[
|v_1|=|v_2|=|v_3| .
\]

Three complex numbers of equal modulus can sum to zero only by forming
an equilateral triangle in the complex plane, hence

\[
\boxed{\Delta\arg=2\pi/3=120^\circ}
\]

follows exactly.

The earlier ``about \(120^\circ\)'' is therefore not a numerical
observation but a theorem from local closure plus equipartition. Since a
\(120^\circ\) configuration contains no pair of opposite signs, this
also explains at once why \(N=4\) has no two-edge zero closures of the
\(N=5\) type.

\subsection{\texorpdfstring{22. The Special Symmetry of
\(N=5\)}{22. The Special Symmetry of N=5}}\label{the-special-symmetry-of-n5}

Among the present series \(N=3\)--16, \(N=5\) is not merely ``a complex
simplex with five vertices and rank 4''; it is \textbf{the clearest
singular case in which phase classification, complex-squared-distance
classification, a three-dimensional readout, and partial zero closures
independent of the generic vertex closures all coincide}.

Here we describe \(N=5\) in detail, including its time evolution, not
only its final state.

We first show overview figures that capture the structure obtained at
\(N=5\) on one page. The detailed subsections below decompose and verify
each element of these figures in turn.

\begin{figure}
\centering
\pandocbounded{\includegraphics[keepaspectratio,alt={N=5: overview of the complex relational distance structure. Five vertices, ten relations, rank 4, and the 3+3+2+2 structure are displayed together.}]{figs/N5_complex_distance_structure_overview_en.png}}
\caption{\(N=5\): overview of the complex relational distance structure.
Five vertices, ten relations, rank 4, and the \(3+3+2+2\) structure are
displayed together.}
\end{figure}

\begin{figure}
\centering
\pandocbounded{\includegraphics[keepaspectratio,alt={N=5: overview of the complete complex-simplex analysis. The two-stage evolution --- rapid expansion and four-class ordering --- is shown on a common time axis.}]{figs/N5_complete_simplex_analysis_overview_en.png}}
\caption{\(N=5\): overview of the complete complex-simplex analysis. The
two-stage evolution --- rapid expansion and four-class ordering --- is
shown on a common time axis.}
\end{figure}

\subsubsection{22.1 Basic geometry: five vertices, ten relations, rank
4}\label{basic-geometry-five-vertices-ten-relations-rank-4}

At \(N=5\) the number of complete pairwise relations is

\[
M=\frac{5\cdot4}{2}=10 .
\]

From the ten complex relational waves \(z_{ij}\) we form the complex
squared distances

\[
d_{ij}^2=z_{ij}^2
\]

and the centered complex Gram matrix

\[
B=-\frac12 JD^2J ,
\]

whose nontrivial rank is

\[
\boxed{\operatorname{rank}B=4} .
\]

The basic structure is thus ``five vertices, ten edges, rank 4'',
corresponding to an ordinary 4-simplex, and this in itself is not
special to \(N=5\). What is special is that \textbf{the phases and
complex squared distances of the ten relations self-organize further
into a small number of highly symmetric classes}.

\subsubsection{\texorpdfstring{22.2 The ten edges finally order into the
four classes
\(3+3+2+2\)}{22.2 The ten edges finally order into the four classes 3+3+2+2}}\label{the-ten-edges-finally-order-into-the-four-classes-3322}

In the metastable state at step 5000 the ten edges converge to the four
classes

\[
A_+:\{12,13,45\},
\]

\[
A_-:\{14,15,23\},
\]

\[
B_+:\{24,35\},
\]

\[
B_-:\{25,34\},
\]

i.e.

\[
\boxed{10=3+3+2+2} .
\]

The important point is that this is not an artificial clustering of
``edges with nearby phases''. Tracking \(a,b,a^2,b^2,ab\) directly for
each edge, edges in the same class converge to the same class even in
the real and imaginary parts of the complex squared distance

\[
z_{ij}^2=(a_{ij}^2-b_{ij}^2)+2i a_{ij}b_{ij} .
\]

\begin{figure}
\centering
\pandocbounded{\includegraphics[keepaspectratio,alt={N=5: final complex-squared-distance classes. The ten edges converge to the four classes 3+3+2+2.}]{figs/N5_final_z2_classes.png}}
\caption{\(N=5\): final complex-squared-distance classes. The ten edges
converge to the four classes \(3+3+2+2\).}
\end{figure}

\subsubsection{\texorpdfstring{22.3 The four classes reduce further to
``two distance directions \(\times\) sign
reversal''}{22.3 The four classes reduce further to ``two distance directions \textbackslash times sign reversal''}}\label{the-four-classes-reduce-further-to-two-distance-directions-times-sign-reversal}

Among the four classes,

\[
A_-=-A_+,
\qquad
B_-=-B_+ .
\]

The four classes are therefore not four independent kinds; essentially
they are

\begin{center}\fbox{\parbox{0.9\linewidth}{\centering $2$ complex squared-distance directions $\times$ sign reversal}}\end{center}

For complex squared distances the sign reversal \(z^2\rightarrow-z^2\)
shifts the phase by \(\pi\), so \(A_+\) and \(A_-\), and \(B_+\) and
\(B_-\), can each be read as the opposite orientations of the same
distance family.

\begin{figure}
\centering
\pandocbounded{\includegraphics[keepaspectratio,alt={N=5: the two distance families and sign reversal. Families A and B, each with its sign reversal, make up the four classes.}]{figs/N5_two_family_components.png}}
\caption{\(N=5\): the two distance families and sign reversal. Families
A and B, each with its sign reversal, make up the four classes.}
\end{figure}

This structure is the basis for reading the \(N=5\) phase classification
not as ``four accidental classes'' but as \textbf{a reversal symmetry
with respect to two basic directions}.

\subsubsection{22.4 Readout as a three-dimensional square
pyramid}\label{readout-as-a-three-dimensional-square-pyramid}

The distance geometry of \(N=5\) itself is a 4-simplex of rank 4. Using
the \(3+3+2+2\) phase/distance classification above, however, the ten
edges can be arranged very naturally as a three-dimensional square
pyramid.

The edges of a square pyramid are

\begin{itemize}
\tightlist
\item
  4 base edges,
\item
  4 lateral edges from the apex to the 4 base vertices,
\item
  2 base diagonals,
\end{itemize}

in total

\[
4+4+2=10 .
\]

All ten relations can therefore be read three-dimensionally as

\[
\boxed{8\text{ outer edges}+2\text{ base diagonals}} .
\]

\begin{figure}
\centering
\pandocbounded{\includegraphics[keepaspectratio,alt={N=5: readout as a three-dimensional square pyramid. The phase/distance classification on the rank-4 complex simplex is visualized three-dimensionally.}]{figs/N5_square_pyramid_interpretation_en.png}}
\caption{\(N=5\): readout as a three-dimensional square pyramid. The
phase/distance classification on the rank-4 complex simplex is
visualized three-dimensionally.}
\end{figure}

This point is important to avoid misreading:

\[
\boxed{
\text{complex simplex rank}=4
}
\]

and

\[
\boxed{
\text{three-dimensional readout of the phase/distance classification}
}
\]

are different observables, and the two do not contradict each other.

\subsubsection{22.5 The four classes do not exist from the beginning of
the
expansion}\label{the-four-classes-do-not-exist-from-the-beginning-of-the-expansion}

Tracing the final \(3+3+2+2\) structure backwards step by step, the
completed pattern is not fixed from the onset of the rapid expansion.

The amplitude outside the parent plane grows rapidly first; afterwards
the phase and distance differences within each class shrink over a long
time.

At \(N=5\), the transverse squared amplitude \(H_\perp\) reaches 99\% of
its final value at about step 449, whereas the error to the final
four-class pattern keeps relaxing until

\[
\le10^{-4}:\ \text{step }2627,
\]

\[
\le10^{-8}:\ \text{step }4923 .
\]

\begin{figure}
\centering
\pandocbounded{\includegraphics[keepaspectratio,alt={N=5: convergence to the four-class structure. The ordering of intra-class phases and complex distances continues after the expansion.}]{figs/N5_four_group_convergence.png}}
\caption{\(N=5\): convergence to the four-class structure. The ordering
of intra-class phases and complex distances continues after the
expansion.}
\end{figure}

The time evolution at \(N=5\) is therefore the two-stage process

\begin{center}\fbox{\parbox{0.9\linewidth}{\centering rapid expansion into new directions $\rightarrow$ ordering into the four-class phase/distance structure}}\end{center}

This observation is the most detailed visualization, at \(N=5\), of the
paper-wide claim that ``the stopping of the expansion and the simplex
symmetrization correspond in time''.

\subsubsection{\texorpdfstring{22.6 The whole structure at
\(N=5\)}{22.6 The whole structure at N=5}}\label{the-whole-structure-at-n5}

In summary, at \(N=5\) several readout levels hold simultaneously for
the same ten relational waves.

\begin{enumerate}
\def\labelenumi{\arabic{enumi}.}
\tightlist
\item
  \textbf{As a complete graph:} \(K_5\), 10 edges.
\item
  \textbf{As distance geometry:} a complex 4-simplex with 5 vertices and
  rank 4.
\item
  \textbf{As a metastable phase classification:} the four classes
  \(3+3+2+2\).
\item
  \textbf{As an algebraic reduction:} two complex distance directions
  and their sign reversals.
\item
  \textbf{As a three-dimensional readout:} 8 outer edges + 2 base
  diagonals of a square pyramid.
\item
  \textbf{As local closure:} the vertex-star zero closure of the 4 edges
  incident to each vertex.
\item
  \textbf{As additional closure:} nontrivial two-edge zero closures not
  explained by the vertex-star span.
\end{enumerate}

\begin{figure}
\centering
\pandocbounded{\includegraphics[keepaspectratio,alt={N=5: complete complex-simplex analysis. Rank, four classes, two distance families, square-pyramid readout, and time evolution are integrated.}]{figs/N5_complete_simplex_analysis_overview_en.png}}
\caption{\(N=5\): complete complex-simplex analysis. Rank, four classes,
two distance families, square-pyramid readout, and time evolution are
integrated.}
\end{figure}

\subsubsection{\texorpdfstring{22.7 Derivation of the nontrivial
two-edge zero closures: 13 and 12 are forced by
\(3+3+2+2\)}{22.7 Derivation of the nontrivial two-edge zero closures: 13 and 12 are forced by 3+3+2+2}}\label{derivation-of-the-nontrivial-two-edge-zero-closures-13-and-12-are-forced-by-3322}

After excluding the trivial 4-edge star closure at each vertex, 13
two-edge zero closures were observed at \(N=5\).

\begin{figure}
\centering
\pandocbounded{\includegraphics[keepaspectratio,alt={N=5: graph of the nontrivial two-edge zero closures. The 13 closures are derived from the 3+3+2+2 sign-paired structure.}]{figs/N5_nontrivial_pair_closure_graph.png}}
\caption{\(N=5\): graph of the nontrivial two-edge zero closures. The 13
closures are derived from the \(3+3+2+2\) sign-paired structure.}
\end{figure}

This is not an independent numerical coincidence. Denote the four
classes by

\[
A_+,\ A_-,\ B_+,\ B_-
\]

with sizes

\[
|A_+|=|A_-|=3,
\qquad
|B_+|=|B_-|=2,
\]

and squared-distance values

\[
q(A_-)=-q(A_+),
\qquad
q(B_-)=-q(B_+) .
\]

Then the two-edge zero closures \(q_e+q_f=0\) are exactly all
combinations between classes of opposite sign, so their number is

\[
\boxed{3\times3+2\times2=13}.
\]

Furthermore, to partition all ten edges into five zero pairs one chooses
independently a bijection between \(A_+\) and \(A_-\) (\(3!\) ways) and
a bijection between \(B_+\) and \(B_-\) (\(2!\) ways), so

\[
\boxed{3!\times2!=12}.
\]

The ``13 closures'' and ``12 exact covers'' found earlier by numerical
search are therefore derived as a \textbf{combinatorial necessity} of
the \(3+3+2+2\) sign-paired structure.

The residual \(1.88\times10^{-10}\) of the best numerical example,

\[
z_{23}^2+z_{45}^2\simeq0,
\]

is now a verification value showing that the numerical implementation
reproduces this derived structure.

\subsubsection{\texorpdfstring{22.8 \(N=5\) is not ``simple because it
is
low-dimensional''}{22.8 N=5 is not ``simple because it is low-dimensional''}}\label{n5-is-not-simple-because-it-is-low-dimensional}

\(N=4\) also has a clear opposite-edge symmetry \(2+2+2\), but its
closures can be explained as linear consequences of the vertex-star
closures. At \(N=5\), by contrast, in addition to the four-class
structure, two-edge closures and exact covers independent of the generic
star span appear.

What makes \(N=5\) special in the present \(N=3\)--16 series is
therefore not merely that it has few edges and is easy to visualize. It
is that

\begin{center}\fbox{\parbox{0.9\linewidth}{\centering simplex symmetry $+$ phase degeneracy $+$ three-dimensional readout $+$ additional zero closures}}\end{center}

coincide at the same \(N\).

\subsection{\texorpdfstring{23. Search for Additional Closures at
\(N=14\)--16}{23. Search for Additional Closures at N=14--16}}\label{search-for-additional-closures-at-n1416}

Over all orders smaller than the trivial \((N-1)\)-edge vertex-star
closure, we searched

\begin{itemize}
\tightlist
\item
  \(N=14\): \(k=2,\ldots,12\),
\item
  \(N=15\): \(k=2,\ldots,13\),
\item
  \(N=16\): \(k=2,\ldots,14\).
\end{itemize}

At \(N=14\) a candidate with

\[
C=9.25\times10^{-7}
\]

exists at \(k=6\), but in the time evolution it stalls near \(10^{-6}\)
and does not converge continuously to \(10^{-8}\), \(10^{-10}\) as at
\(N=5\).

It is therefore classified as a quasi-closure rather than a strict
additional closure.

At \(N=15,16\) no additional closure clearly below \(10^{-6}\) was found
at any order searched.

\begin{center}\rule{0.5\linewidth}{0.5pt}\end{center}

\subsection{24. Reconstruction of the Axiom System in This Paper ---
What Is Settled and What Is
Not}\label{reconstruction-of-the-axiom-system-in-this-paper-what-is-settled-and-what-is-not}

Previously the two basic axioms were

\[
\sum z_m^2=0,
\qquad
U^n=I .
\]

From the self-consistent real antisymmetric eigenmode we have now
organized analytically

\begin{center}\fbox{\parbox{0.9\linewidth}{\centering self-consistency $\Rightarrow$ complex rotating pair $\Rightarrow
\sum z_m^2=0
\Rightarrow
S^1$ compact phase orbit}}\end{center}

However,

\[
S^1\text{ compactness}\not\Rightarrow U^n=I .
\]

Finite recurrence requires, in addition, rational locking of the phase
increment.

The axiom-reduction conclusion of this paper is therefore revised to:
\textbf{zero-square closure and the complex rotating structure can be
derived from self-consistency, but the full derivation of \(U^n=I\)
requires an additional rational phase-locking mechanism}.

\subsection{25. Relation to Prior Work}\label{relation-to-prior-work}

\subsubsection{25.1 Chronology of the
research}\label{chronology-of-the-research}

The numerical model of this paper was not built for the purpose of
implementing or reproducing the external prior work listed below.

The chronological order of the research is

\begin{center}\fbox{\parbox{0.9\linewidth}{\centering independent numerical experiments $\rightarrow$ discovery of the phenomenon $\rightarrow$ later literature survey $\rightarrow$ recognition of structural relations}}\end{center}

The external literature is therefore cited not as the starting
assumptions of the model, but to locate independently obtained results
in the existing research space.

\subsubsection{25.2 Complex structure from real
structure}\label{complex-structure-from-real-structure}

Aste {[}E1{]} discusses the possibility that a complex structure emerges
from the linear dynamics of a finite-dimensional real Euclidean Hilbert
space.

The common concern with the present paper is that ``complex numbers need
not be a primitive input from the outset''.

The difference is that here the self-consistent fixed point of a real
antisymmetric action built from complete pairwise relations generates
the complex rotating pair.

\subsubsection{25.3 Low-dimensional expanding spacetime from
high-dimensional degrees of
freedom}\label{low-dimensional-expanding-spacetime-from-high-dimensional-degrees-of-freedom}

In the Lorentzian type IIB / IKKT matrix model, the possibility that a
low-dimensional expanding spacetime emerges dynamically from many matrix
degrees of freedom has been studied numerically {[}E2,E3{]}.

The present paper shares the structural feature that rank generation and
direction selection arise from many relational degrees of freedom.

The action, the degrees of freedom, the probability measure, and the
numerical methods differ, however, and this paper is not a reproduction
of the IKKT model.

\subsubsection{25.4 Inflation and dynamical
compactification}\label{inflation-and-dynamical-compactification}

Kihara et al.~{[}E4{]} study, in a ten-dimensional Einstein--Yang--Mills
system, exponential expansion on the four-dimensional side and
compactification of the extra dimensions within one dynamical system.

In the seedless system of this paper all principal axes expand
isotropically; in future harmonic-seed experiments we plan to test an
anisotropic separation in which some directions become macroscopic while
others remain compact.

In this sense that work is an important object of future comparison.

\begin{center}\rule{0.5\linewidth}{0.5pt}\end{center}

\subsection{26. Discussion --- the Onset Mechanism Is Quantified as the
Linear Instability of a Relative
Equilibrium}\label{discussion-the-onset-mechanism-is-quantified-as-the-linear-instability-of-a-relative-equilibrium}

\subsubsection{26.1 Seedless onset and the fixed-point
residual}\label{seedless-onset-and-the-fixed-point-residual}

With no explicit seed added and \(Z_0=v\), we varied the requested
tolerance of \texttt{make\_parent} over
\(10^{-6},10^{-8},10^{-10},10^{-12}\). The explanatory variable used in
the regression, however, is not the requested tolerance itself but the
final fixed-point residual \(\varepsilon_{\rm parent}\) actually
measured in each run.

\begin{figure}
\centering
\pandocbounded{\includegraphics[keepaspectratio,alt={N=5: sweep of the make\_parent residual tolerance. A smaller residual delays the onset while the exponential growth rate is almost unchanged.}]{figs/N5_tol_sweep_timeseries.png}}
\caption{\(N=5\): sweep of the \protect\texttt{make\_parent} residual
tolerance. A smaller residual delays the onset while the exponential
growth rate is almost unchanged.}
\end{figure}

The measured values are as follows.

{\def\LTcaptype{none} % do not increment counter
\begin{longtable}[]{@{}
  >{\raggedleft\arraybackslash}p{(\linewidth - 6\tabcolsep) * \real{0.2500}}
  >{\raggedleft\arraybackslash}p{(\linewidth - 6\tabcolsep) * \real{0.2500}}
  >{\raggedleft\arraybackslash}p{(\linewidth - 6\tabcolsep) * \real{0.2500}}
  >{\raggedleft\arraybackslash}p{(\linewidth - 6\tabcolsep) * \real{0.2500}}@{}}
\toprule\noalign{}
\begin{minipage}[b]{\linewidth}\raggedleft
requested tol
\end{minipage} & \begin{minipage}[b]{\linewidth}\raggedleft
measured parent residual \(\varepsilon_{\rm parent}\)
\end{minipage} & \begin{minipage}[b]{\linewidth}\raggedleft
onset step \(f\ge10^{-8}\)
\end{minipage} & \begin{minipage}[b]{\linewidth}\raggedleft
growth rate /step
\end{minipage} \\
\midrule\noalign{}
\endhead
\bottomrule\noalign{}
\endlastfoot
\(10^{-6}\) & \(3.8728081613\times10^{-7}\) & 72 & 0.17483734 \\
\(10^{-8}\) & \(1.8154355031\times10^{-9}\) & 134 & 0.17252539 \\
\(10^{-10}\) & \(5.0849521984\times10^{-11}\) & 176 & 0.17251311 \\
\(10^{-12}\) & \(2.3845913288\times10^{-13}\) & 238 & 0.17251360 \\
\end{longtable}
}

The logarithms of the residuals are not equally spaced. The onset step
must therefore not be regressed against the requested tolerances
\(10^{-6},10^{-8},10^{-10},10^{-12}\) alone. Using the \textbf{measured
fixed-point residuals} in the table,

\begin{figure}
\centering
\pandocbounded{\includegraphics[keepaspectratio,alt={N=5: onset versus the measured fixed-point residual. t\_\{\textbackslash rm onset\} is linear in -\textbackslash ln(\textbackslash text\{residual\}) to high precision.}]{figs/N5_tol_sweep_onset_log_residual.png}}
\caption{\(N=5\): onset versus the measured fixed-point residual.
\(t_{\rm onset}\) is linear in \(-\ln(\text{residual})\) to high
precision.}
\end{figure}

the regression

\[
t_{\rm onset}=a[-\ln\varepsilon_{\rm parent}]+b
\]

gives

\[
a=11.616225,
\qquad b=-99.563139,
\qquad R^2=0.99999201.
\]

The predicted onsets for the four points are 71.940, 134.236, 175.766,
238.058 steps, differing from the integer detected values by at most
0.236 step. The regression statistics of the previous version are
therefore consistent with the raw data; the apparent contradiction arose
from identifying the ``requested tolerance'' with the ``measured
fixed-point residual''.

This result shows that the intrinsic floor of the seedless onset is the
\texttt{make\_parent} fixed-point residual: the residual controls the
onset time, while in the regime of sufficiently small residuals the
dominant growth rate converges to about 0.17251/step.

\subsubsection{26.2 The rotating-frame Jacobian explains rank-4 and the
growth rate
simultaneously}\label{the-rotating-frame-jacobian-explains-rank-4-and-the-growth-rate-simultaneously}

We linearized numerically the 20-real-dimensional Jacobian of the
rotating-frame map, in which the one-step global phase rotation of the
parent state has been removed.

\begin{figure}
\centering
\pandocbounded{\includegraphics[keepaspectratio,alt={N=5: rotating-frame Floquet spectrum. It shows the gap between the dominant unstable real doubly degenerate eigenvalue and the second one. All unstable multipliers lie on the real axis.}]{figs/N5_floquet_spectrum.png}}
\caption{\(N=5\): rotating-frame Floquet spectrum. It shows the gap
between the dominant unstable real doubly degenerate eigenvalue and the
second one. All unstable multipliers lie on the real axis.}
\end{figure}

The largest unstable multiplier is real (zero imaginary part) and
appears twice, with value

\[
|\mu_1|=1.090086569 .
\]

In the measured spectrum (\texttt{floquet\_spectrum.csv}) both the
dominant multiplier \(\mu_1\) and the weaker second unstable multiplier
\(\mu_2=1.052603212\) appear as real doubly degenerate eigenvalues, with
a clear gap in growth rate between them. What can be established
directly from the data is therefore that the dominant unstable
eigenspace is real two-dimensional and gives the fastest unstable
directions.

All twenty multipliers show a clear pair structure. Global phase
equivariance alone, \[
F(e^{i\alpha}Z)=e^{i\alpha}F(Z),
\] does not automatically imply, for an arbitrary perturbation
\(\delta Z\), \[
DF(Z)[i\,\delta Z]=i\,DF(Z)[\delta Z],
\] i.e.~that the whole Jacobian is complex linear. To attribute the
symmetry origin of the real doubling to the global phase symmetry, one
would have to show separately that the Jacobian commutes with the
complex structure. In this paper the doubling itself is treated as an
observed fact of the numerical spectrum, and its symmetry origin is left
as a future analytical problem.

The sum of the log-moduli of all twenty multipliers is about
\(-0.0855\), which is negative. Since the rotating-frame Jacobian is the
derivative of the 20-real-dimensional state map, \[
\sum_i\ln|\mu_i|=\ln|\det J|<0
\] shows that local state volume contracts near the parent relative
equilibrium. This does not contradict the exact conservation of the
total norm, and it is consistent with the possibility of interpreting
the simplex ordering as attractor-like. Restricting the contraction to
particular ``angular directions'', however, would require a separate
tangent-space decomposition.

Hence the selection rule

\[
\boxed{\text{parent rank-2}+\text{dominant unstable 2D}=\text{rank-4}}
\]

is obtained.

Furthermore, from the amplitude exponent

\[
\lambda_A=\ln\mu_1=0.0862571143
\]

the squared-amplitude exponent

\[
2\lambda_A=0.172514229
\]

is predicted, which agrees with the growth rate of \(H_\perp\) in the
direct time evolution, about 0.172513.

Moreover,

\[
\frac1{\lambda_A}=11.5932
\]

agrees with the slope 11.6162 of onset versus residual obtained from the
tolerance sweep.

For the onset of the rapid expansion at \(N=5\), therefore, the
quantitative mechanism

\begin{center}\fbox{\parbox{0.9\linewidth}{\centering self-consistent relative equilibrium $\rightarrow$ linear instability $\rightarrow$ dominant 2D unstable eigenspace $\rightarrow$ rank-4 and $H_\perp$ exponential growth}}\end{center}

holds.

The second unstable real doubly degenerate eigenvalue

\[
|\mu_2|=1.052603212
\]

also exists and predicts a slower secondary growth exponent
\(2\ln\mu_2\simeq0.10253\).

\begin{figure}
\centering
\pandocbounded{\includegraphics[keepaspectratio,alt={N=5: finite-difference stability of the Floquet multipliers. The dominant multipliers are stable against the finite-difference step.}]{figs/N5_floquet_fd_stability.png}}
\caption{\(N=5\): finite-difference stability of the Floquet
multipliers. The dominant multipliers are stable against the
finite-difference step.}
\end{figure}

\subsubsection{26.3 Structural similarity to preheating-type mode
transfer}\label{structural-similarity-to-preheating-type-mode-transfer}

Because the total norm is conserved exactly in this system, the rapid
expansion is not an energy generation but an exponential transfer from
the coherent parent mode to the orthogonal modes. Structurally it is
closer to parametric mode amplification / preheating than to
cosmological inflation proper.

This paper does not claim identity with the standard preheating
equations, however. We only point out the dynamical similarity: Floquet
instability, depletion of the parent component, and long-time
redistribution after saturation.

Moreover, even at step 5000, \(H_\parallel=0.322005\) and
\(H_\perp=0.677995\): the parent-plane component has not vanished.
Taking an uncorrelated isotropic state on 20 real dimensions as a naive
reference, the expected projection ratio onto a fixed two-dimensional
parent plane is \(2/20=0.1\), so the observed value 0.322 is about 3.22
times that. Since we have not shown that the metastable distribution of
this map is uniform on the 20-dimensional sphere, this is treated as ``a
diagnostic suggesting residual memory of the initial parent plane'', and
its strict statistical significance is to be tested in future seed
series.

\subsubsection{\texorpdfstring{26.4 Seed sweep of the \(N=5\)
moduli}{26.4 Seed sweep of the N=5 moduli}}\label{seed-sweep-of-the-n5-moduli}

Re-running 5000 steps for 8 parent random seeds, the final class sizes
were

\[
3+3+2+2
\]

in every run, and both distance-family moduli converged to about 0.1.

The relative phase between the two distance families, taken modulo sign,
varied from run to run.

\begin{figure}
\centering
\pandocbounded{\includegraphics[keepaspectratio,alt={N=5: relative phase between the two distance families across seeds. 3+3+2+2 is reproduced, while the relative phase is not fixed uniquely.}]{figs/N5_moduli_relative_phase_seed_sweep.png}}
\caption{\(N=5\): relative phase between the two distance families
across seeds. \(3+3+2+2\) is reproduced, while the relative phase is not
fixed uniquely.}
\end{figure}

This supports the possibility that \(3+3+2+2\) and equimodularity are
selected rigidly, whereas a flat direction remains in the relative phase
between the distance families. Longer runs with more seeds are needed.

\subsection{27. Limitations and Future
Work}\label{limitations-and-future-work}

The additional analyses of this paper have narrowed some of the earlier
open problems.

\begin{enumerate}
\def\labelenumi{\arabic{enumi}.}
\tightlist
\item
  \textbf{Onset of the rapid expansion}: at \(N=5\) it has been
  quantified as the linear instability of a relative equilibrium.
  Whether the same Floquet structure holds at other \(N\) is to be
  checked.
\item
  \textbf{Rank-4 selection}: at \(N=5\) it is explained by the
  two-real-dimensional dominant unstable space. The real-time growth of
  the second unstable pair is to be separated directly.
\item
  \textbf{Stopping}: the reason for the absence of unbounded growth is
  resolved by total-norm conservation. The remaining problem is why the
  system orders into an equimodular null simplex.
\item
  \textbf{Thermalization / ordering}: the spectral entropy eventually
  becomes maximal but not monotonically. The \(N\)-dependence of the
  relaxation time and the presence or absence of FPUT-type
  quasi-integrability are to be examined.
\item
  \textbf{Moduli at \(N=5\)}: whether the relative phase is a true flat
  direction or locks over very long times is to be confirmed with many
  seeds and long runs.
\item
  \textbf{\(U^n=I\)}: a compact phase orbit has been derived from
  self-consistency, but finite recurrence requires a rational
  phase-locking mechanism.
\item
  \textbf{General classification of the \(N=5\) singularity}: the
  conditions for the existence of opposite-sign phase classes are to be
  classified as balanced edge-class designs on \(K_N\).
\item
  \textbf{Floquet spectra at other \(N\)}: the correspondence between
  the dimension of the unstable space and the simplex ordering is to be
  organized as a series over \(N=3\)--16.
\item
  \textbf{Parent-plane memory}: at \(N=5\), \(H_\parallel=0.322\)
  remains even at step 5000. The seed dependence, \(N\) dependence, and
  long-time limit of this residual fraction are to be examined.
\item
  \textbf{Symmetry origin of the real doubly degenerate eigenvalues}:
  the reason why all twenty multipliers of the rotating-frame Jacobian
  come in pairs and the unstable multipliers appear as real doubles was
  treated as a numerical observation. Since global phase equivariance
  alone does not imply complex linearity of the Jacobian, commutation
  with the complex structure or another symmetry must be shown
  analytically.
\end{enumerate}

\subsection{28. Conclusion}\label{conclusion}

This paper re-audited the rapid expansion of the seedless closed
relational-wave system and promoted some of the numerical observations
to exact theorems and part of the onset mechanism to a linear-stability
problem.

The main results are as follows.

\begin{enumerate}
\def\labelenumi{\arabic{enumi}.}
\tightlist
\item
  The Cayley update of a real antisymmetric generator is real orthogonal
  and conserves \(Z^\dagger Z\) and \(Z^TZ\) exactly.
\item
  Hence the growth of \(H_\perp\) is not an external injection but an
  internal transfer from \(H_\parallel\), and
  \(H_\perp\le H_{\rm total}\) gives the kinematic upper bound of the
  expansion.
\item
  The \(K/\sigma_{\max}\) normalization corresponds exactly, at the same
  state, to a rescaling of the Cayley step. The roughly 6.8\%
  growth-rate difference observed at finite step is predicted as
  \(g_{\rm pred}\simeq1.0600\) from the first-order step-convergence law
  \(g(n)=1.1596346-4.10498/n\) of §12 and \(n_{\rm eff}\simeq41.2\) of
  the raw-\(K\) run, agreeing with the measured \(1.05874\) to about
  0.13\%. Both updates therefore converge to the same continuum limit,
  and the possibility that normalization or step size artificially
  generated the expansion is excluded.
\item
  In the seedless \(N=5\) system, reducing the \texttt{make\_parent}
  fixed-point residual delays the onset in proportion to
  \(-\ln\varepsilon\) while the growth rate is unchanged.
\item
  The rotating-frame Jacobian has a dominant unstable real doubly
  degenerate eigenvalue \(\mu_1=1.090086569\), faster than the second
  real doubly degenerate eigenvalue \(\mu_2=1.052603212\). The numerical
  spectrum establishes that the dominant unstable space is real
  two-dimensional, and adding it to the rank-2 parent explains the
  initial selection of rank-4. The symmetry origin of the real doubling
  requires a separate proof of commutation between the Jacobian and the
  complex structure.
\item
  \(2\ln|\mu_1|=0.172514\) agrees with the directly measured \(H_\perp\)
  growth rate of about 0.172513, and \(1/\ln|\mu_1|=11.593\) agrees with
  the onset-regression slope 11.616 against the \textbf{measured
  fixed-point residual}. This triple consistency is the central
  quantitative result on the \(N=5\) onset mechanism.
\item
  Zero-square closure of every vertex star is equivalent, in the
  centered representation, to all simplex vertices lying on the complex
  null cone.
\item
  Together with equipartition, the metastable state is described
  uniformly as an \textbf{equimodular null complex simplex}.
\item
  The \(120^\circ\) phase difference at \(N=4\) is derived exactly from
  local closure plus equipartition.
\item
  The 13 nontrivial two-edge closures and 12 exact covers at \(N=5\) are
  derived exactly from the \(3+3+2+2\) sign-paired structure as
  \(3\times3+2\times2=13\) and \(3!2!=12\).
\item
  In the 8-seed sweep at \(N=5\), \(3+3+2+2\) and equimodularity were
  reproduced in every run, whereas the relative phase between the two
  distance families was not fixed uniquely, indicating a possible
  modulus.
\item
  From self-consistency one can derive the complex rotating structure,
  zero-square closure, and a compact \(S^1\) phase orbit, but \(U^n=I\)
  requires separate rational phase locking.
\end{enumerate}

The onset of the rapid expansion at \(N=5\) can therefore be described,
at least within numerical linearization, as one dynamical sequence,

\begin{center}\fbox{\parbox{0.9\linewidth}{\centering self-consistent rank-2 relative equilibrium $\rightarrow$ dominant 2D linear instability $\rightarrow$ rank-4 $\rightarrow$ exponential internal mode transfer $\rightarrow$ bounded saturation $\rightarrow$ equimodular null-simplex ordering}}\end{center}

\subsection{Location of Figure Files}\label{location-of-figure-files}

The figures are numbered consecutively from Figure 1 in order of
appearance.

The figures are stored in the repository under
\texttt{dimensional-generation/} \texttt{electron-recoil/}
\texttt{seed-removal-metastable/} \texttt{chatgpt-followup/} (repository
path, romanized) and referenced by relative paths from there (the Zenodo
release bundles the full set of figures in
\path{nbody_self_consistent_inflation_figures_v1.zip}). The figures used
in the additional experiments of this paper comprise the normalization
comparison, the transfer from the parent plane to the orthogonal
directions, the spectral entropy, the tolerance sweep, the
onset--residual law, the Floquet spectrum, the finite-difference
stability, and the \(N=5\) moduli seed sweep. Figures 7, 8, 11, and 13
are the English versions of the corresponding Japanese infographics,
regenerated from the same analysis data.

\subsection{Self-Citations}\label{self-citations}

\textbf{{[}K1{]}} N. Kihara, \emph{Anonymous Equal-Amplitude
Composite-Wave Model: Basic Axiom System v9 --- Pure Definition}, 2026.

\textbf{{[}K2{]}} N. Kihara, ``Linear Upper Bound on Generator Rank and
Three-Direction Saturation in N-Body Complete Pairwise Relational
Waves'', 2026. Concept DOI: 10.5281/zenodo.21465898.

\textbf{{[}K3{]}} N. Kihara, ``Plane-Decomposition Readout in N-Body
Fixed-Generator Systems'', 2026. Concept DOI: 10.5281/zenodo.21468959.

\textbf{{[}K4{]}} N. Kihara, ``Onset of Spontaneous State Splitting and
the Threefold Classification of Its Outcomes in Closed N-Body
Relational-Wave Systems'', 2026. Concept DOI: 10.5281/zenodo.21486233.

\textbf{{[}K5{]}} N. Kihara, ``The Number of Waves Is the Resolution of
the System'', 2026. Concept DOI: 10.5281/zenodo.21486544.

\textbf{{[}K6{]}} N. Kihara, ``The Three-Direction Structure Arising
from the Antisymmetric Generator of Complete Pairwise Relational
Waves'', 2026.

\textbf{{[}K7{]}} N. Kihara, ``Emergence of Three-Direction Space in
Closed N-Body Relational-Wave Systems'', 2026.

\textbf{{[}K8{]}} N. Kihara, ``Causal Separation of the Time Structure
of Three-Direction Formation in Closed N-Body Relational Wave Systems
via Two-Stage Seed Removal'', 2026. Version DOI:
10.5281/zenodo.21614403; Concept DOI: 10.5281/zenodo.21614402.\\
note: https://note.com/kiharanoriaki/n/n48a02cd70f47

\subsection{External Prior Work}\label{external-prior-work}

\textbf{{[}E1{]}} A. Aste, ``Origin of the Complex Structure of Quantum
Mechanics,'' arXiv:1905.12894, 2019.\\
https://arxiv.org/abs/1905.12894

\textbf{{[}E2{]}} T. Aoki, M. Hirasawa, Y. Ito, J. Nishimura, A.
Tsuchiya, ``On the structure of the emergent 3D expanding space in the
Lorentzian type IIB matrix model,'' \emph{Progress of Theoretical and
Experimental Physics} 2019, 093B03 (2019). DOI: 10.1093/ptep/ptz092.

\textbf{{[}E3{]}} M. Hirasawa, K. N. Anagnostopoulos, T. Azuma, K.
Hatakeyama, J. Nishimura, S. Papadoudis, A. Tsuchiya, ``The effects of
SUSY on the emergent spacetime in the Lorentzian type IIB matrix
model,'' arXiv:2407.03491, 2024.\\
https://arxiv.org/abs/2407.03491

\textbf{{[}E4{]}} H. Kihara, M. Nitta, M. Sasaki, C.-M. Yoo, I. Zaballa,
``Dynamical Compactification and Inflation in Einstein-Yang-Mills Theory
with Higher Derivative Coupling,'' \emph{Physical Review D} \textbf{80},
066004 (2009). DOI: 10.1103/PhysRevD.80.066004.

\begin{center}\rule{0.5\linewidth}{0.5pt}\end{center}

\subsection{Note on Reproducibility}\label{note-on-reproducibility}

The re-run data, programs, and figure generation used in this paper are
stored in the following verification packages created on 2026-08-26 (the
\(\gamma\) continuum bundle on 2026-08-25).

\begin{itemize}
\item
  \path{K_sigma_normalization_artifact_test_N4_N5_20260826.zip}
\item
  \path{N5_gamma_continuum_test_bundle_20260825.zip}
\item
  \path{N3_N4_complex_simplex_complete_analysis_20260826.zip}
\item
  \path{N5_complex_simplex_complete_analysis_20260826.zip}
\item
  \path{N6_N7_complex_simplex_complete_analysis_20260826.zip}
\item
  \path{N8_N9_complex_simplex_complete_analysis_20260826.zip}
\item
  \path{N10_N11_complex_simplex_complete_analysis_20260826.zip}
\item
  \path{N12_N13_complex_simplex_complete_analysis_20260826.zip}
\item
  \path{N14_N15_complex_simplex_complete_analysis_20260826.zip}
\item
  \path{N16_complex_simplex_complete_analysis_20260826.zip}
\item
  \path{N3_N16_partial_zero_closure_analysis_20260826.zip}
\item
  \path{N3_N16_nontrivial_zero_closure_analysis_20260826.zip}
\item
  \path{N14_N16_complete_nontrivial_zero_closure_search_20260826.zip}
\item
  \path{N5_dynamics_followup_theorems_and_stability_20260826.zip}
\end{itemize}

With these programs, raw data, and figure scripts, the main numerical
results of this paper can be reproduced.

\end{document}
